\section{Lifting torus regions to the plane}%
\label{sec:peps_torus_region_lifts}

This section lifts a region of the torus to a union of closed unit squares in
the plane. For a simply connected region the lift has no holes, and walks in
the complement can be chosen with prescribed crossing numbers of the two
seams.

\begin{lemma}[Integer fibers and affine lifts of the additive circle]%
    \label{thm:peps_additive_circle_lift_coordinates}
    \lean{TNLean.PEPS.exists_intCast_eq_of_addCircle_eq_zmod_val,
        TNLean.PEPS.eq_affine_of_addCircle_eq}
    \leanok
    Let $n>0$ be an integer, let $z\in\mathbb Z/n\mathbb Z$, and let
    $\bar z\in\{0,\ldots,n-1\}$ be its standard representative.
    If $x\in\mathbb R$ has the same image as $\bar z$ in
    $\mathbb R/n\mathbb Z$, there is an integer $j$ with $j=x$ and
    $j\equiv z\pmod n$.

    For any real period $p$, let $f:[0,1]\to\mathbb R$ be continuous.
    If its image in the additive circle satisfies $[f(t)]=[f(0)]+[td]$
    for every $t\in[0,1]$, then $f(t)=f(0)+td$ throughout the interval.
    These are auxiliary
    arithmetic and covering-space identities used for the cell and
    edge paths of the block domains in \cite[\S6.3]{Schuch2010PEPS}.
\end{lemma}
\begin{proof}\leanok
    The first equality means that $x-\bar z$ is an integer multiple
    of $n$, so $x$ is an integer with the stated residue. In the
    second assertion, $f$ and the affine path have identical projections
    and agree initially; uniqueness of a covering lift on the connected
    interval makes them equal.
\end{proof}

\begin{lemma}[Finite compact cell images]%
    \label{thm:peps_finite_compact_cell_images}
    \lean{TNLean.PEPS.locallyPathConnectedSpace_iUnion_range}
    \leanok
    Let $K$ be compact and locally path connected, let $Y$ be Hausdorff,
    and let $f_i:K\to Y$ be a finite family of continuous maps. Then
    $\bigcup_i f_i(K)$ is locally path connected.
\end{lemma}
\begin{proof}\leanok
    The finite disjoint union of copies of $K$ is compact and locally
    path connected. Its natural continuous surjection onto the stated
    union is closed, hence is a quotient map. Local path connectedness
    passes through this quotient map.
\end{proof}

\begin{lemma}[Closed-cell regions and geometric coordinate lifts]%
    \label{thm:peps_torus_region_realization}
    \lean{TNLean.PEPS.torusRealProjection_apply,
        TNLean.PEPS.torusRealProjection,
        TNLean.PEPS.continuous_torusRealProjection,
        TNLean.PEPS.torusVertexPoint,
        TNLean.PEPS.torusVertexPoint_eq_projection,
        TNLean.PEPS.torusVertexPoint_intCast,
        TNLean.PEPS.torusVertexPoint_right,
        TNLean.PEPS.torusVertexPoint_up,
        TNLean.PEPS.torusClosedUnitCell,
        TNLean.PEPS.torusRegionRealization,
        TNLean.PEPS.torusVertexPoint_mem_torusRegionRealization,
        TNLean.PEPS.torusRightUnitPath,
        TNLean.PEPS.torusUpUnitPath,
        TNLean.PEPS.torusRightUnitPath_apply,
        TNLean.PEPS.torusUpUnitPath_apply,
        TNLean.PEPS.isCompact_torusClosedUnitCell,
        TNLean.PEPS.isCompact_torusRegionRealization,
        TNLean.PEPS.torusRegionRealization_locallyPathConnectedSpace,
        TNLean.PEPS.torusRightUnitPath_mem_torusRegionRealization,
        TNLean.PEPS.torusUpUnitPath_mem_torusRegionRealization,
        TNLean.PEPS.exists_isTorusRegionIntegerLift_of_continuousLift,
        TNLean.PEPS.exists_continuousLift_torusRegionRealization,
        TNLean.PEPS.exists_isTorusRegionIntegerLift_of_isSimplyConnected}
    \leanok
    Let $w,H>0$, let
    $\pi:\mathbb R^2\to(\mathbb R/w\mathbb Z)\times
        (\mathbb R/H\mathbb Z)$ be the quotient map, and let $p_v$
    be the image of the native vertex coordinates. Define the closed
    cell and the realized region by
    \begin{align}
        Q_v & =p_v+\pi([-1/2,1/2]^2), &
        X_R & =\bigcup_{v\in R}Q_v.
    \end{align}
    The cells and their finite union are compact, and $X_R$ is locally
    path connected. Every native site in $R$ lies in $X_R$. The paths
    $t\mapsto p_v+\pi(t,0)$ and $t\mapsto p_v+\pi(0,t)$ have the
    native neighboring sites as their endpoints; whenever both sites
    belong to $R$, the whole path lies in $X_R$, including across a
    periodic seam.

    Any continuous lift $F:X_R\to\mathbb R^2$ of the inclusion
    yields integer coordinates $L:R\to\mathbb Z^2$ which project to
    the native vertices and preserve every internal horizontal and
    vertical unit step. In particular, if $X_R$ is simply connected
    and a site $o\in R$ is fixed, such a continuous lift and such
    integer coordinates exist.

    This closed-cell convention is an auxiliary geometric realization
    of the block domains drawn in \cite[\S6.3]{Schuch2010PEPS}.
    Simple connectedness is the ordinary topological hypothesis on
    the cell union; no disk predicate, group flatness assumption,
    or entropy conclusion is introduced. Its relation to arbitrary
    disk regions in Theorem~6.9 remains recorded in
    \cite{gap:rmp_peps_quantum_double_g_isometry}.
\end{lemma}
\begin{proof}\leanok
    \uses{thm:peps_additive_circle_lift_coordinates, thm:peps_finite_compact_cell_images}
    The quotient projection is continuous, so each projected closed
    square is compact. The finite coproduct of the compact convex
    squares maps continuously and surjectively onto $X_R$. This is a
    closed quotient map into a Hausdorff space, and local path
    connectedness therefore passes to $X_R$. Splitting a native unit
    path at its midpoint places its first half in the initial cell and
    its second half in the final cell.

    Applying covering-space lifting to the two coordinate maps into
    $\mathbb R/w\mathbb Z$ and $\mathbb R/H\mathbb Z$ separately
    assigns integer coordinates to every lift of a native site. Along an
    internal unit path, uniqueness of the lift identifies the lifted
    path with the corresponding affine unit segment. Its endpoints
    thus satisfy the exact integer unit-step equations.
\end{proof}

\begin{lemma}[The planar lift of the closed-cell region]%
    \label{thm:peps_torus_region_planar_lift}
    \lean{TNLean.PEPS.continuousLift_torusClosedUnitCell,
        TNLean.PEPS.continuousLift_torusRegionRealization_isClosedEmbedding,
        TNLean.PEPS.torusRegionPlanarRealization,
        TNLean.PEPS.range_continuousLift_torusRegionRealization,
        TNLean.PEPS.torusRegionLiftHomeomorph,
        TNLean.PEPS.exists_integerLift_torusRegionPlanarRealization}
    \leanok
    \uses{thm:peps_torus_region_realization}
    Let $F:X_R\to\mathbb R^2$ be a continuous lift of the inclusion
    into the real torus. On every native cell,
    \begin{align}
        F(p_v+\pi(d)) & =F(p_v)+d,
                      & d          & \in[-1/2,1/2]^2.
    \end{align}
    The lift is a closed embedding. There are integer site coordinates
    $L(v)=F(p_v)$ which project to the original native vertices and
    preserve every internal unit step, and
    \begin{align}
        F(X_R) & =P_L:=\bigcup_{v\in R}
        \bigl(L(v)+[-1/2,1/2]^2\bigr).
    \end{align}
    Thus $F$ induces a homeomorphism from the torus cell union
    to this finite planar cell union.

    This is an auxiliary geometric consequence of the covering lift,
    supporting the block domains of \cite[\S6.3]{Schuch2010PEPS}.
    The homeomorphism transfers simple connectedness to the planar
    union, which is the domain of the exterior contour argument below.
\end{lemma}
\begin{proof}\leanok
    \uses{thm:peps_additive_circle_lift_coordinates}
    The radial segment from $p_v$ to $p_v+\pi(d)$ stays in the same
    closed cell. Uniqueness of its covering lift identifies the lifted
    segment with $t\mapsto F(p_v)+td$. Two neighboring cells have a
    common side midpoint; evaluating its lift from the two cells shows
    that their lifted centers differ by the prescribed unit step.
    The fibers over native sites have integer coordinates. These
    identities identify the image cell by cell. Since $\pi\circ F$
    is the inclusion, $F$ is injective. Compactness of $X_R$ and the
    Hausdorff property of the plane make this continuous injection a
    closed embedding, giving the stated homeomorphism.
\end{proof}

\begin{lemma}[Separation of distinct period copies of a lifted region]%
    \label{thm:peps_torus_region_period_separation}
    \lean{TNLean.PEPS.torusRealProjection_periodVector,
        TNLean.PEPS.torusPeriodVector,
        TNLean.PEPS.torusPeriodVector_eq_zero_iff,
        TNLean.PEPS.torusRealProjection_eq_iff_exists_periodVector,
        TNLean.PEPS.continuousLift_torusRegionRealization_periodTranslate_disjoint,
        TNLean.PEPS.torusRegionPlanarRealization_periodTranslate_disjoint,
        TNLean.PEPS.preimage_torusRegionRealization_eq_iUnion_periodTranslate,
        TNLean.PEPS.torusRegionPlanarRealization_cell_periodTranslate_disjoint,
        TNLean.PEPS.exists_integerLift_periodTranslate_disjoint_of_isSimplyConnected}
    \leanok
    \uses{thm:peps_torus_region_realization,
        thm:peps_torus_region_planar_lift}
    Let the torus periods $w,H$ be positive, and let
    $F:X_R\to\mathbb R^2$ be a continuous lift of the inclusion.
    For $j=(j_1,j_2)\in\mathbb Z^2$, write
    $\delta_j=(j_1w,j_2H)$. Two plane points have the same projection
    exactly when their difference is some $\delta_j$, and
    \begin{align}
        \pi^{-1}(X_R) & =\bigcup_{j\in\mathbb Z^2}(P_L+\delta_j).
    \end{align}
    If $j\ne0$, then
    \begin{align}
        F(X_R)\cap\bigl(F(X_R)+\delta_j\bigr) & =\varnothing, \\
        P_L\cap(P_L+\delta_j)                 & =\varnothing.
    \end{align}
    In particular, individual closed cells in distinct period copies
    cannot meet, even along a side or at a corner. If the cell
    union $X_R$ is simply connected and a site in $R$ is fixed, there
    exists an integer unit-step lift $L$ satisfying this separation
    for every nonzero $j$.

    This is an auxiliary geometric statement for the block domains of
    \cite[\S6.3]{Schuch2010PEPS}. The exterior path-existence argument
    remains separate; see \cite{gap:rmp_peps_quantum_double_g_isometry}.
\end{lemma}
\begin{proof}\leanok
    The additive-circle fibers are integer multiples of the two periods.
    Comparing a point with the lift of its projection gives the stated
    union of all period copies. The projection of $\delta_j$ is zero. If
    $F(q)+\delta_j=F(p)$, projecting to the torus gives $q=p$,
    because the lift is a section of the projection. Subtracting
    $F(p)$ then gives $\delta_j=0$. Positivity of both periods
    forces $j=0$, a contradiction. The exact image description in
    Lemma~\ref{thm:peps_torus_region_planar_lift} gives the statement
    for $P_L$, and restricting to individual cells gives the last
    separation assertion. Simple connectedness yields the continuous
    lift and hence the required integer coordinates.
\end{proof}

% Geometric follow-up target (SCP10, local source lines 1935--1990):
% Let \Gamma be the native torus graph with both dimensions at least three,
% let \Gamma[R] be nonempty and connected, and suppose X_R is simply connected.
% For every crossing edge f write a_f\in R and b_f\in R^c for its endpoints.
% Given f_0,f and an internal walk P:a_{f_0}\leadsto a_f, prove that there is
% an exterior walk Q:b_{f_0}\leadsto b_f in \Gamma[R^c] such that
% \operatorname{wind}(\iota Q)=
% \operatorname{wind}([b_{f_0},a_{f_0}]*(\iota P)*[a_f,b_f]).
% Here \iota includes an induced-graph walk into \Gamma and * concatenates
% walks in traversal order. Appending the chosen exterior root paths on both
% sides then gives a complement loop with the winding of the combined
% boundary route, whose holonomy is m_f m_{f_0}^{-1}.

\begin{lemma}[No finite holes in simply connected closed-cell regions]%
    \label{thm:peps_integer_cell_no_holes}
    \lean{TNLean.PEPS.sum_integerGrid_curl_eq_boundary,
        TNLean.PEPS.integerClosedCellUnion,
        TNLean.PEPS.integerClosedCellUnion_locallyPathConnectedSpace,
        TNLean.PEPS.not_finite_integerCellHole_of_isSimplyConnected,
        TNLean.PEPS.infinite_integerCellComplement_component_of_isSimplyConnected,
        TNLean.PEPS.integerClosedCellUnion_image_eq_torusRegionPlanarRealization,
        TNLean.PEPS.exists_integerLift_no_finite_cellHole_of_isSimplyConnected}
    \leanok
    \uses{thm:peps_torus_region_realization, thm:peps_torus_region_planar_lift}
    For a finite set $A\subset\mathbb Z^2$, the union
    \begin{align}
        P_A & =\bigcup_{a\in A}\bigl(a+[-1/2,1/2]^2\bigr)
        \label{eq:peps_integer_cell_union}
    \end{align}
    is locally path connected. For any real edge values $H_a,V_a$ and finite $K\subset\mathbb Z^2$,
    \begin{align}
        \sum_{a\in K}\bigl(H_a+V_{a+(1,0)}-H_{a+(0,1)}-V_a\bigr)
        \label{eq:peps_integer_grid_curl}
    \end{align}
    equals the signed sum over the exposed boundary edges of $K$. If $P_A$ is simply connected, no finite $K\subset\mathbb Z^2$ with $K\setminus A\ne\varnothing$ satisfies $a+\{\pm e_1,\pm e_2\}\subset K\cup A$ for all $a\in K$; hence every $\ell^1$-connected component of $\mathbb Z^2\setminus A$ is infinite. Consequently, if $X_R$ is simply connected and $o\in R$, there exists a unit-step integer lift $L:R\to\mathbb Z^2$ (Lemma~\ref{thm:peps_torus_region_planar_lift}) such that every missing-cell component of $P_L=P_{L(R)}$ is infinite (an auxiliary no-hole step for \cite[Theorem~6.9]{Schuch2010PEPS}).
\end{lemma}
\begin{proof}\leanok
    \uses{thm:peps_finite_compact_cell_images}
    Local path connectedness follows from Lemma~\ref{thm:peps_finite_compact_cell_images}. Suppose a finite $K\subset\mathbb Z^2$ satisfies $a+\{\pm e_1,\pm e_2\}\subset K\cup A$ for all $a\in K$ and contains $0\notin A$. Simple connectedness of $P_A$ yields a continuous angular lift $\theta:P_A\to\mathbb R$; subtracting $\theta/(2\pi)$ from the principal argument defines a corner potential whose horizontal increments vanish and whose vertical increment equals $\mathbf 1_{\{\text{edge crosses }(-\infty,0)\times\{0\}\}}$. Along $\partial K\subset P_A$, (\ref{eq:peps_integer_grid_curl}) vanishes by interior cancellation yet equals $-1$ from the branch cut, a contradiction. Lemma~\ref{thm:peps_torus_region_planar_lift} transfers this to $X_R\cong P_L$.
\end{proof}

\begin{definition}[Safe graph of an exterior integer-cell collar]%
    \label{def:peps_integer_exterior_collar_graph}
    \lean{TNLean.PEPS.integerCellCenter,
        TNLean.PEPS.integerClosedCell,
        TNLean.PEPS.integerExteriorCollar,
        TNLean.PEPS.IsIntegerCellNear,
        TNLean.PEPS.integerExteriorBand,
        TNLean.PEPS.integerExteriorCollarGraph}
    \leanok
    \uses{thm:peps_integer_cell_no_holes}
    Let $A\subset\mathbb Z^2$ be finite, with occupied closed-cell union
    $P_A$. Its exterior collar is
    \begin{align}
        E_A=\bigl(P_A+[-3/4,3/4]^2\bigr)\setminus P_A.
    \end{align}
    The exterior band consists of the unoccupied integer centers $q$
    such that $\|q-a\|_\infty\leq1$ for some $a\in A$.
    Join two such centers by an edge precisely when they are horizontal
    or vertical neighbors and one common occupied center $a\in A$
    satisfies both $\|q-a\|_\infty\leq1$ and $\|r-a\|_\infty\leq1$.
    The common-center requirement controls the whole segment, rather
    than its endpoints alone.
\end{definition}

\begin{lemma}[Integer centers of the exterior collar]%
    \label{thm:peps_integer_exterior_collar_centers}
    \lean{TNLean.PEPS.mem_integerClosedCell_iff,
        TNLean.PEPS.mem_integerClosedCell_add_of_abs_le_one,
        TNLean.PEPS.integerCellCenter_mem_closedCellUnion_iff,
        TNLean.PEPS.integerCellCenter_mem_exteriorCollar_iff}
    \leanok
    \uses{def:peps_integer_exterior_collar_graph}
    The closed cell at $a$ consists exactly of the points whose two
    coordinate distances from $a$ are at most $1/2$. An integer center
    belongs to $P_A$ exactly when it belongs to $A$, and belongs to
    $E_A$ exactly when it belongs to the exterior band.
\end{lemma}
\begin{proof}\leanok
    A distinct integer center cannot be within $1/2$ in both coordinates.
    In the expanded cells the coordinate distances are at most $5/4$;
    integer distances are therefore at most one. Conversely, a distance
    at most one splits into an offset of at most $1/2$ within the cell
    and an offset of at most $1/2$ within the collar square.
\end{proof}

\begin{lemma}[Continuous collar paths from safe exterior walks]%
    \label{thm:peps_integer_exterior_collar_paths}
    \lean{TNLean.PEPS.integerExteriorCollarGraph_segment_subset,
        TNLean.PEPS.joinedIn_integerExteriorCollar_of_reachable}
    \leanok
    \uses{def:peps_integer_exterior_collar_graph}
    Every edge of the safe exterior graph has its entire straight
    segment inside $E_A$. Hence any graph walk between two band centers
    gives a continuous path between them contained in $E_A$.
    This is an auxiliary geometric assertion for
    \cite[proof of Theorem~6.9]{Schuch2010PEPS}. The contour theorem below derives connectivity of the safe graph
    from simple connectedness of the occupied closed-cell union.
\end{lemma}
\begin{proof}\leanok
    \uses{thm:peps_integer_exterior_collar_centers}
    Both endpoint centers lie within coordinate distance one of their
    common occupied center. Convexity preserves these bounds along the
    segment, which is therefore inside the collar enlargement. If a
    closed occupied cell met a horizontal or vertical unit segment,
    integer arithmetic would force its center to be one of the two
    endpoints. Both endpoints are unoccupied, so the segment is outside
    $P_A$. Concatenating the edge segments proves the path assertion.
\end{proof}

\begin{definition}[Occupied square contact graph]%
    \label{def:peps_integer_cell_contact_graph}
    \lean{TNLean.PEPS.integerCellContactGraph}
    \leanok
    Let $A\subset\mathbb Z^2$ be finite, and let $Q_a$ be the closed
    unit square centered at $a$. The contact graph has vertex set $A$;
    distinct $a,b$ are adjacent when $|a_1-b_1|\leq1$ and
    $|a_2-b_2|\leq1$. Thus shared corners, as well as shared sides,
    are retained.
\end{definition}

\begin{lemma}[Connectivity through occupied square contacts]%
    \label{thm:peps_integer_cell_contact_connectivity}
    \lean{TNLean.PEPS.isIntegerCellNear_of_integerClosedCell_inter_nonempty,
        TNLean.PEPS.integerCellContactGraph_preconnected_of_isPreconnected,
        TNLean.PEPS.integerCellContactGraph_connected_of_isSimplyConnected}
    \leanok
    \uses{def:peps_integer_cell_contact_graph}
    Write $P_A=\bigcup_{a\in A}Q_a$. If $P_A$ is preconnected, then
    every two occupied centers can be joined in the contact graph.
    In particular, if $P_A$ is simply connected, its contact graph
    is connected. This auxiliary geometric consequence supports the
    block decomposition in \cite[Theorem~6.9]{Schuch2010PEPS}, local
    source lines 1935--1990. Four-neighbor connectivity of the
    occupied centers is not assumed.
\end{lemma}

\begin{proof}\leanok
    A point in $Q_a\cap Q_b$ places each coordinate of $a$ and $b$
    within one half of that point, so their coordinate differences
    have absolute value at most one. Fix an occupied center $u$ and
    partition $A$ into the centers reachable from $u$ and the
    remaining centers. The corresponding square unions are closed,
    since $A$ is finite. If both are nonempty, preconnectedness of
    $P_A$ forces an intersection between the unions. The intersecting
    cells give a contact edge across the partition, a contradiction.
    Simple connectedness also gives nonemptiness of $P_A$, hence of $A$.
\end{proof}


\begin{definition}[Horizontal-ray potential of a finite edge flow]%
    \label{def:peps_integer_contour_ray_potential}
    \lean{TNLean.PEPS.integerContourRayPotential}\leanok
    Index each oriented lattice edge by its starting corner
    $a-(1/2,1/2)$, with horizontal edges pointing east and vertical
    edges pointing north. For a finitely supported signed vertical
    flow $V:\mathbb Z^2\to\mathbb Z$, define its ray potential at the
    cell center $q$ by
    \begin{align}
        W(q)=\sum_{\substack{a_2=q_2 \\q_1<a_1}}V(a).
    \end{align}
    The sum is finite because $V$ is finitely supported.
\end{definition}

\begin{lemma}[Discrete gradients of the contour ray potential]%
    \label{thm:peps_integer_contour_ray_potential}
    \lean{TNLean.PEPS.integerContourRayPotential_east,
        TNLean.PEPS.integerContourRayPotential_north}
    \leanok
    \uses{def:peps_integer_contour_ray_potential}
    Write $e=(1,0)$ and $n=(0,1)$. For any finitely supported signed
    vertical flow,
    \begin{align}
        W(q+e)-W(q)=-V(q+e).
    \end{align}
    If a finitely supported horizontal flow $H$ satisfies conservation
    \begin{align}
        H(a)+V(a)-H(a-e)-V(a-n)=0
        \qquad(a\in\mathbb Z^2),
    \end{align}
    then also $W(q+n)-W(q)=H(q+n)$.
    These are auxiliary winding-potential identities for the contour
    argument associated with \cite[proof of Theorem~6.9]{Schuch2010PEPS}.
    Identification of the flow with an exposed contour and
    connectivity of that contour are separate geometric assertions.
\end{lemma}
\begin{proof}\leanok
    Moving the ray one cell east removes exactly the vertical edge
    indexed by $q+e$. Conservation writes the difference between the
    vertical flow and its north translate as the east translate of
    the horizontal flow minus the horizontal flow itself. Applying the
    additive finite ray sum to this equality, translating the
    finite supports, and using the east-increment identity telescopes the
    horizontal difference to the single term $H(q+n)$.
\end{proof}

\begin{lemma}[Connected missing-cell lattice of a simply connected region]%
    \label{thm:peps_integer_cell_exterior_connected}
    \lean{TNLean.PEPS.integerCellComplement_connected_of_isSimplyConnected,
        TNLean.PEPS.exists_integerLift_integerCellComplement_connected_of_isSimplyConnected}
    \leanok
    \uses{thm:peps_integer_cell_no_holes}
    Let $A\subset\mathbb Z^2$ be finite, and suppose its closed-cell union
    $P_A=\bigcup_{a\in A}(a+[-1/2,1/2]^2)$ is simply connected. Then
    the four-neighbor lattice on $\mathbb Z^2\setminus A$ is connected.
    Consequently, a simply connected torus region admits an
    integer unit-step lift whose entire missing-cell lattice is connected.

    These are plane-complement walks. They need not stay in the
    exterior collar, so their torus projections are not yet asserted
    to avoid all period copies of the region. The following contour theorem restricts the walks to the genuine
    exterior collar, as required in
    \cite[proof of Theorem~6.9]{Schuch2010PEPS}.
\end{lemma}
\begin{proof}\leanok
    Enclosing $A$ in a finite integer box allows all vertices outside that box
    to be connected to one exterior corner by integer row and column
    walks avoiding the box. Every missing-cell component is infinite
    by the preceding no-hole theorem, and hence contains a vertex
    outside the box. Each component therefore meets the same exterior
    component, proving connectedness. The torus lift given
    by the no-hole theorem has the same infinitude property, so the
    argument applies to its integer centers as well.
\end{proof}

\begin{lemma}[Connected exterior cell contour]%
    \label{thm:peps_integer_boundary_contour_connected}
    \lean{TNLean.PEPS.integerExteriorCollarGraph_connected_of_isSimplyConnected}
    \leanok
    \uses{def:peps_integer_exterior_collar_graph}
    Let $A\subset\mathbb Z^2$ be finite and suppose that the union
    $P_A=\bigcup_{a\in A}(a+[-1/2,1/2]^2)$ is simply connected.
    Let $B_A$ consist of the unoccupied centers whose coordinate distances
    from some occupied center are at most one. Join four-neighbor centers
    in $B_A$ whenever one occupied center is within coordinate distance
    one of both. Then this graph is connected.
    This is auxiliary exterior-collar geometry for
    \cite[proof of Theorem~6.9]{Schuch2010PEPS}, local source lines
    1935--2076. Shared corners of occupied cells are allowed; simple
    connectedness refers to the closed-cell union $P_A$.
\end{lemma}

\begin{proof}\leanok
    \uses{thm:peps_integer_cell_exterior_connected,
        thm:peps_integer_cell_contact_connectivity,
        thm:peps_integer_contour_ray_potential}
    Orient each exposed cell edge with its occupied cell on the left.
    At an edge endpoint, continue with a right turn when the diagonal
    cell is occupied, straight ahead when only the forward cell is
    occupied, and a left turn otherwise. The reverse rule makes this
    successor a permutation of the finite exposed-edge set.

    Fix one successor orbit and assign weight one to its edges and zero
    to the remaining edges. The resulting finite flow is conserved.
    Its horizontal-ray winding potential is constant on missing cells,
    since their four-neighbor graph is connected, and vanishes far from
    the occupied cells. It is also constant on occupied cells: the edge
    identities give equality across shared sides, and the right-turn
    rule gives equality across shared corners. The occupied contact
    graph is connected. Across any exposed edge the winding difference
    is its weight. Thus every exposed edge has the same weight, and all
    belong to the chosen successor orbit.

    Consecutive contour edges have exterior centers that are equal or
    joined by one or two four-neighbor steps with a common occupied-cell
    support. Every missing center in $B_A$ can reach one of these
    exterior centers using the same local support. These walks prove
    graph connectivity. Nonemptiness follows by taking an occupied
    center with maximal first coordinate and the missing center to its east.
\end{proof}


\begin{lemma}[Exclusion of foreign period copies from an exterior collar]%
    \label{thm:peps_torus_region_exterior_collar}
    \lean{TNLean.PEPS.disjoint_torusRegionPlanarRealization_add_collar_periodTranslate,
        TNLean.PEPS.torusRealProjection_not_mem_of_mem_planarExteriorCollar,
        TNLean.PEPS.exists_integerLift_exteriorCollar_of_isSimplyConnected}
    \leanok
    \uses{thm:peps_torus_region_period_separation,
        thm:peps_torus_region_realization,
        thm:peps_torus_region_planar_lift}
    Let $w,H>0$, and let a continuous lift of the
    closed-cell region $X_R$ have planar image
    $P=\bigcup_{v\in R}(L(v)+[-1/2,1/2]^2)$ with integer centers.
    Put $E=[-3/4,3/4]^2$, and let $P_j=P+(wj_1,Hj_2)$ be its period
    copies. Then
    \begin{align}
        (P+E)\cap P_j                   & =\varnothing\qquad(j\neq0), \\
        \pi\bigl((P+E)\setminus P\bigr) & \subseteq X_R^c.
    \end{align}
    Here addition of sets is Minkowski addition. If $X_R$ is simply
    connected and contains a native site, there is an integer unit-step
    lift $L$ with this exterior-collar property; no covering lift is
    assumed as input in that conclusion.

    This is auxiliary exterior geometry for
    \cite[proof of Theorem~6.9]{Schuch2010PEPS}. Collar paths between exterior integer centers follow from the
    connected safe exterior graph established above.
\end{lemma}
\begin{proof}\leanok
    Disjoint closed unit squares with integer centers differ by at
    least two in some coordinate; otherwise their midpoint belongs
    to both. A point in an enlarged cell and in a foreign unit cell
    would put the two center coordinates within
    $1/2+3/4+1/2=7/4<2$ in both directions, a contradiction.
    The full preimage of $X_R$ is the union of the period copies $P_j$.
    A point of the exterior collar lies in neither $P$ nor any foreign
    copy, proving the projection assertion. In the simply connected
    case, covering-space lifting gives the genuine lift and its
    integer centers, so the same argument applies.
\end{proof}

\begin{theorem}[Simply connected integer-cell lift with an exterior collar]%
    \label{thm:peps_torus_integer_cell_lift}
    \lean{TNLean.PEPS.exists_integerLift_simplyConnected_exteriorCollar_of_isSimplyConnected}
    \leanok
    Let $R$ be a finite region of a torus with positive integer
    periods, and assume its closed-cell realization $X_R$
    is simply connected. There exists an integer lift $L$ of the
    occupied vertices, preserving each internal unit step. Put
    $A=L(R)$ and let $P_A$ be the union of the closed unit squares
    centered at $A$. Then $P_A$ is simply connected, and every point
    of its exterior collar projects outside $X_R$:
    \begin{align}
        \pi\!\left(\bigl(P_A+[-\tfrac34,\tfrac34]^2\bigr)\setminus P_A\right)
        \cap X_R & =\varnothing.
        \label{eq:peps_integer_lift_collar}
    \end{align}
    Both conclusions hold for the same integer lift. This is an
    auxiliary geometric consequence for the block decomposition in
    \cite[Theorem~6.9]{Schuch2010PEPS}, local source lines 1935--1990.
\end{theorem}

\begin{proof}\leanok
    \uses{thm:peps_torus_region_planar_lift,
        thm:peps_torus_region_exterior_collar}
    Simple connectedness gives a continuous lift of the covering
    projection over $X_R$. Its image is homeomorphic to $X_R$ and is
    precisely $P_A$, so $P_A$ is simply connected. Distinct period
    copies of the lifted cells are disjoint. The separation of
    integer centers implies that the enlargement by
    $[-\tfrac34,\tfrac34]^2$ misses every other period copy.
    A point in the exterior collar also misses $P_A$ itself; hence it
    cannot project to $X_R$, proving (\ref{eq:peps_integer_lift_collar}).
\end{proof}


\begin{lemma}[Closure gauges from a given integer lift]%
    \label{thm:peps_torus_region_integer_lift_gauge}
    \lean{TNLean.PEPS.IsTorusRegionIntegerLift.injective,
        TNLean.PEPS.IsTorusRegionIntegerLift,
        TNLean.PEPS.int_ediv_step_eq_add_seam_indicator,
        TNLean.PEPS.torusRegionLiftGauge,
        TNLean.PEPS.torusRegionLiftGauge_right,
        TNLean.PEPS.torusRegionLiftGauge_up,
        TNLean.PEPS.torusRegionLiftGauge_gradient,
        TNLean.PEPS.regularRegionGaugeResidual_torusRegionLiftGauge,
        TNLean.PEPS.torusRegionLiftRootGauge,
        TNLean.PEPS.torusRegionLiftRootGauge_gradient,
        TNLean.PEPS.regularRegionTreeGauge_eq_torusRegionLiftRootGauge,
        TNLean.PEPS.regularRegionTreeCycleResidual_torusClosure_eq_one}
    \leanok
    \uses{thm:peps_regular_region_tree_gauge,
        thm:peps_regular_torus_closure_cut_identification}
    Let both torus circumferences $w,H$ be at least three, and let
    $R$ be a set of native torus vertices. Suppose a map
    $L:R\to\mathbb Z^2$ projects modulo $(w,H)$ to the original
    vertices and preserves every internal rightward and upward unit
    step. In particular, $L$ is injective. The arithmetic deck coordinate
    $\lfloor j/n\rfloor$ increases by one under $j\mapsto j+1$ precisely
    when $j+1\equiv0\pmod n$, and otherwise stays unchanged. For commuting $g,h\in G$,
    define
    \begin{align}
        k_v & =h^{\lfloor L_x(v)/w\rfloor}
        g^{-\lfloor L_y(v)/H\rfloor}.
    \end{align}
    On a native rightward step the gauge changes by left multiplication
    by $h$ precisely at the horizontal seam; on an upward step it
    changes by left multiplication by $g^{-1}$ precisely at the vertical
    seam. For every internal edge, ordered from $s$ to $t$,
    \begin{align}
        k_tk_s^{-1} & =u_e,\qquad k_t^{-1}u_ek_s=1,
    \end{align}
    where the native closure assignment is $h^{-1}$ on a horizontal
    seam edge, $g$ on a vertical seam edge, and the identity elsewhere.
    Right normalization $k_v\mapsto k_vk_o^{-1}$ preserves the gradient.
    For finite $G$, this normalized gauge equals the gauge obtained from
    any spanning tree rooted at $o$; hence all its cycle residuals equal
    the identity.

    This is an auxiliary consequence of the given arithmetic lift,
    supporting the seam and disentangling arguments in
    \cite[\S5.3 and Theorem~6.9]{Schuch2010PEPS}. It neither defines a
    disk nor establishes lifts for arbitrary disk regions. The latter
    geometric theorem and the removal of sector-dependent crossing
    factors remain separate, as recorded in
    \cite{gap:rmp_peps_quantum_double_g_isometry}.
\end{lemma}
\begin{proof}\leanok
    Euclidean division shows that the quotient of an integer coordinate
    increases by one exactly when its residue wraps through zero.
    The unit-step hypotheses therefore give the two gauge-step
    identities; commutativity moves the vertical factor past the
    horizontal factor. A seam reverses the native ordered-edge
    orientation, yielding $h^{-1}$ and $g$ in the respective gradients.
    All other internal gradients are the identity. Right normalization
    cancels in each gradient. Uniqueness in the tree-gradient bijection
    identifies the normalized gauge with the spanning-tree gauge.
\end{proof}

\begin{lemma}[Crossing numbers of projected integer steps]%
    \label{thm:peps_torus_integer_unit_step_winding}
    \lean{TNLean.PEPS.torusDirectedWinding_eq_of_integerUnitStep,
        TNLean.PEPS.torusIntegerDeckCoordinate}
    \leanok
    \uses{thm:peps_torus_walk_winding,
        thm:peps_torus_region_integer_lift_gauge}
    Let the torus periods $w,H$ be at least three. For $a\in\mathbb Z^2$,
    put
    \begin{align}
        D(a) & =\left(\left\lfloor a_1/w\right\rfloor,
        \left\lfloor a_2/H\right\rfloor\right).
    \end{align}
    If $a,b$ differ by one horizontal or vertical unit step, then the
    oriented crossing pair of their projected native edge is $D(b)-D(a)$.
    This applies to interior, boundary, and exterior steps alike.
\end{lemma}
\begin{proof}\leanok
    In a positive horizontal step, Euclidean division changes precisely
    when the standard residue crosses the horizontal seam. The vertical
    case is identical. Reversing a step negates its crossing pair.
\end{proof}

\begin{lemma}[Exterior lattice walks from a given plane path]%
    \label{thm:peps_torus_exterior_path_approximation}
    \lean{TNLean.PEPS.exists_complementWalk_of_exterior_planePath}
    \leanok
    \uses{thm:peps_torus_region_realization,
        thm:peps_torus_walk_winding}
    Let the torus periods $w,H$ be at least three, and let $X_R$ be
    the closed-cell realization of a finite native region $R$.
    Let $a,b\in\mathbb Z^2$ project to exterior native vertices $v,t$.
    Suppose a continuous plane path $\gamma$ joins $a$ to $b$ and
    its torus projection avoids $X_R$ at every time. There is a walk
    $Q:v\to t$ in the induced native complement with crossing numbers
    \begin{align}
        N(Q)=\left(
        \left\lfloor\frac{b_1}{w}\right\rfloor
        -\left\lfloor\frac{a_1}{w}\right\rfloor,
        \left\lfloor\frac{b_2}{H}\right\rfloor
        -\left\lfloor\frac{a_2}{H}\right\rfloor
        \right).
    \end{align}
    Thus the lattice walk has the prescribed lifted endpoint displacement.
    This is an auxiliary path approximation for the complement argument
    in \cite[proof of Theorem~6.9]{Schuch2010PEPS}. The path $\gamma$ is taken as given in this approximation theorem;
    the preceding contour construction provides such paths for the
    required exterior centers of simply connected regions.
\end{lemma}
\begin{proof}\leanok
    An exterior point belongs to an open rectangle containing one cell,
    two cells sharing a side, or four cells sharing a corner, all with
    exterior lattice centers. At a side or corner every touching closed
    cell must be exterior, since the point avoids $X_R$. These rectangles
    cover the lifted exterior. Subdividing the compact parameter interval
    places each subpath in one rectangle, where rounded endpoints
    can be connected through exterior lattice centers
    by a horizontal and a vertical step. Concatenating gives an integer
    lattice walk between $a$ and $b$. Projection gives a native
    complement walk. The crossing pair of each unit step is the change
    in its integer deck coordinates; summation telescopes to the stated
    endpoint difference.
\end{proof}

\begin{lemma}[Lifted exterior endpoints of a region route]%
    \label{thm:peps_torus_region_boundary_lift}
    \lean{TNLean.PEPS.torusVertexPoint_mem_torusRegionRealization_iff,
        TNLean.PEPS.exists_exterior_integerUnitStep_of_isTorusRegionIntegerLift,
        TNLean.PEPS.exists_exterior_integerEndpoints_of_regionWalk}
    \leanok
    \uses{thm:peps_torus_integer_unit_step_winding, thm:peps_torus_region_planar_lift}
    A native site center belongs to the closed-cell union $X_R$ exactly
    when its site belongs to $R$. For the crossing and walk assertions,
    take both torus periods at least three. Suppose $L$ is an integer lift
    preserving the unit steps within $R$, and put
    $P=\bigcup_{v\in R}(L(v)+[-1/2,1/2]^2)$.
    An edge from $v\in R$ to $t\notin R$ lifts to a unit step from
    $L(v)$ to an integer point $b$ projecting to $t$, with
    \begin{align}
        b        & \in (P+[-3/4,3/4]^2)\setminus P, \\
        N([v,t]) & =D(b)-D(L(v)).
    \end{align}
    More generally, let $p:v_0\to v$ be any native region walk, and
    let $t_0,t$ be exterior neighbors of its endpoints. There are lifted
    exterior integer endpoints $b_0,b$ in the same collar, projecting
    to $t_0,t$, such that
    \begin{align}
        N([t_0,v_0]*p*[v,t]) & =D(b)-D(b_0).
    \end{align}
    Concatenation is written in traversal order. This gives the
    endpoints and winding pair required in the complement argument of
    \cite[proof of Theorem~6.9]{Schuch2010PEPS}; the connected safe exterior graph gives a collar path between them.
\end{lemma}
\begin{proof}\leanok
    \uses{thm:peps_region_walk_winding}
    A site center in a closed cell differs from its integer center,
    modulo a period vector, by an integer in each interval $[-1/2,1/2]$.
    Both differences vanish, proving the center assertion. Lifting a boundary
    edge by its horizontal or vertical unit direction places its exterior
    endpoint at the sum of its midpoint in the interior closed cell and
    a half-unit displacement, so it belongs to the collar. It cannot
    belong to $P$, since projection would put its exterior native center
    in $X_R$. The preceding theorem gives the crossing pair of this step.
    The interior walk has crossing pair $D(L(v))-D(L(v_0))$; adding the
    two boundary crossings cancels the interior endpoint coordinates.
\end{proof}

\begin{lemma}[Exterior replacement with prescribed winding]%
    \label{thm:peps_torus_complement_path_replacement}
    \lean{TNLean.PEPS.exists_complementWalk_of_integerExteriorCollarGraph_connected,
        TNLean.PEPS.exists_complementBoundaryLoops_of_integerExteriorCollarGraph_connected}
    \leanok
    \uses{def:peps_integer_exterior_collar_graph}
    Let both torus periods be at least three. Suppose a region has an
    integer unit-step lift, its safe exterior graph is connected, and
    its exterior collar projects outside the torus region.
    Every route which enters through one crossing bond, follows an
    interior path, and leaves through another crossing bond has a
    complementary replacement with the same two seam crossing numbers.

    In the auxiliary geometric construction for
    \cite[proof of Theorem~6.9]{Schuch2010PEPS}, fix in particular rooted
    spanning trees in the region and its complement and a reference crossing
    bond. For each crossing bond there is a loop based at the complementary
    root, wholly within the induced complement, with the winding of the
    corresponding two-tree boundary route, chosen independently of all
    group-valued closure labels.
\end{lemma}
\begin{proof}\leanok
    \uses{thm:peps_integer_exterior_collar_paths,
        thm:peps_torus_region_boundary_lift,
        thm:peps_torus_exterior_path_approximation}
    Lifting the two exterior crossing endpoints places them in the
    integer band, with deck-coordinate difference equal to the
    crossing numbers of the interior route. A safe graph walk gives
    a continuous collar path between them, and projection exclusion together with
    lattice-path approximation gives a native complementary walk with
    that same deck-coordinate difference. Prepending the complementary
    tree path from the root to the first crossing and appending the
    tree path from the second crossing back to the root yields the loop, since crossing
    numbers add under concatenation.
\end{proof}

\begin{lemma}[Complementary boundary loops of a simply connected region]%
    \label{thm:peps_torus_simply_connected_boundary_loops}
    \lean{TNLean.PEPS.exists_complementBoundaryLoops_of_isSimplyConnected}\leanok
    In the exterior geometric construction of
    \cite[proof of Theorem~6.9]{Schuch2010PEPS}, let both torus periods be at
    least three, and suppose the closed-cell realization of $R$ is simply
    connected. Choose rooted spanning trees in $R$ and its induced complement
    and fix a reference crossing bond. There is a complementary root loop for
    every crossing bond with the same winding as its two-tree boundary route,
    chosen before all group-valued closure labels without any additional
    exterior-path, collar-connectivity, or winding-identity premise.
\end{lemma}
\begin{proof}\leanok
    \uses{thm:peps_torus_complement_path_replacement,
        thm:peps_integer_boundary_contour_connected,
        thm:peps_torus_integer_cell_lift}
    Choosing one covering lift whose planar integer-cell union is simply
    connected and whose exterior collar projects outside the torus
    region gives a safe exterior graph that is connected by the boundary contour
    theorem, so the preceding replacement construction applies to the
    chosen two-tree routes.
\end{proof}


\begin{definition}[The rectangular grid obtained by cutting the torus seams]%
    \label{def:peps_torus_rectangular_grid}
    \lean{TNLean.PEPS.torusNonseamGraph,
        TNLean.PEPS.torusNonseamGraph_le_torusGraph}
    \leanok
    \uses{def:peps_torus_geometry_translation}
    On the torus vertices, retain only consecutive coordinate values
    without wraparound. The result is the ordinary rectangular grid,
    a subgraph of the periodic square lattice.
\end{definition}

\begin{lemma}[Connectivity of rectangles and their complements]%
    \label{thm:peps_torus_rectangle_cut_connectivity}
    \lean{TNLean.PEPS.torusNonseamGraph_compl_rectangle_connected,
        TNLean.PEPS.torusGraph_compl_rectangle_connected,
        TNLean.PEPS.torusNonseamGraph_rectangle_connected,
        TNLean.PEPS.torusGraph_rectangle_connected,
        TNLean.PEPS.nonempty_boundaryEdge_torusRectangle}
    \leanok
    \uses{def:peps_torus_rectangular_grid}
    A bounded coordinate rectangle with positive side lengths is connected
    in the ordinary grid, hence also in the periodic lattice.
    If both side lengths are strictly smaller than the torus dimensions,
    its complement is connected even without the wrap bonds. This last
    assertion also holds for coordinate intervals clipped by the bounds
    of the grid. A nonempty bounded rectangle starting at a positive
    horizontal coordinate has a crossing bond on its left side.
\end{lemma}

\begin{proof}\leanok
    Inside the rectangle, consecutive coordinate changes connect any two vertices.
    For the complement, choosing one row and one column missed by the
    rectangle lets every outside vertex reach their intersection along
    an outside row or column, while a vertex at the lower left corner and
    its immediate left neighbour give the crossing bond.
\end{proof}

\begin{lemma}[Connected complement of a contiguous simply connected block]%
    \label{thm:peps_torus_simply_connected_complement}
    \lean{TNLean.PEPS.torusGraph_compl_connected_of_isSimplyConnected,
        TNLean.PEPS.nonempty_torusRegionBoundaryEdge_of_isSimplyConnected}
    \leanok
    Let both torus circumferences be at least three. Let $R$ be a finite
    set of sites whose induced nearest-neighbour graph is connected and
    whose closed-cell realization is simply connected. Then the
    induced graph on the complementary sites is connected, and at least
    one bond crosses between $R$ and its complement.

    Consequently, spanning trees and roots on both sides of the cut,
    and a numbering of the crossing bonds by $n+1$ labels, may be chosen
    without further hypotheses. Occupied graph connectivity is distinct
    from simple connectedness of the closed-cell union: cells meeting
    only at a corner need not be connected by lattice bonds. This is the
    contiguous-block geometry used in Section~6.3 of
    \cite{Schuch2010PEPS}, rather than a claim that every simply connected
    closed-cell subset is a contiguous block.
\end{lemma}

\begin{proof}\leanok
    \uses{thm:peps_torus_integer_cell_lift,
        thm:peps_integer_exterior_collar_paths,
        thm:peps_integer_boundary_contour_connected,
        thm:peps_torus_complement_path_replacement,
        thm:peps_torus_rectangle_cut_connectivity}
    Lift the occupied region to a simply connected finite integer-cell
    union whose exterior collar projects outside the region. The exposed
    contour joins its nearby missing centres within the safe exterior
    graph. These joins yield complementary lattice walks between all
    outside endpoints of crossing bonds. Every complementary site has
    a walk to an occupied site in the full torus graph; its first crossing
    joins it, within the complement, to such an endpoint. Hence all
    complementary sites are connected. The collar also provides a
    complementary site, so a walk from an occupied site crosses at least
    one boundary bond.
\end{proof}


\begin{lemma}[Closure-region Grams from an integer lift]%
    \label{thm:peps_torus_region_integer_lift_gram}
    \lean{TNLean.PEPS.IsGIsometric.regularTwistedOpenRegionMatrix_torusRegionLift,
        TNLean.PEPS.IsGIsometric.exists_positive_gram_torusClosureRegionMatrix_of_integerLift,
        TNLean.PEPS.IsGIsometric.rank_torusClosureRegionMatrix_of_integerLift}
    \leanok
    Let the torus circumferences be at least three, and let $A$ be a
    native regular $G$-isometric four-leg tensor. Equip a region $R$ with
    an integer lift preserving all internal unit steps, and let $g,h$
    commute. For the region matrix $M_R^{g,h}$ of the native
    closure, the explicit lift gauge $k$ gives
    \begin{align}
        M_R^{g,h}(\sigma,x) & =M_R(\sigma,F_{g,h}(x)).
    \end{align}
    If $R$ is connected, there is a positive scalar $c_R$ such that
    \begin{align}
        (M_R^{g,h})^\dagger M_R^{g,h}
         & =c_R\Pi_{\partial R}^{k,u(g,h)}.
    \end{align}
    When the boundary has $b>0$ bonds,
    $\operatorname{rank}M_R^{g,h}=|G|^{b-1}$.
    These conclusions follow from the given arithmetic lift, without
    any separate flatness or block Gram hypothesis. They support the
    seam-deformation argument of \cite[Theorem~6.9]{Schuch2010PEPS};
    they do not establish lifts for arbitrary disk regions or identify
    boundary maps for superpositions of distinct closures
    \cite{gap:rmp_peps_quantum_double_g_isometry}.
\end{lemma}

\begin{proof}\leanok
    \uses{thm:peps_torus_region_integer_lift_gauge, thm:peps_regular_region_gauge_gram, thm:peps_regular_torus_closure_cut_identification}
    Expressing the native four-leg tensor on its four incident graph edges
    preserves local regular isometry. Since the arithmetic lift gauge
    has the native closure assignment as its internal gradient,
    every internal residual equals the identity, and applying
    Lemma~\ref{thm:peps_regular_region_gauge_gram} to the resulting
    region matrix with the boundary map of this explicit gauge yields the claims.
\end{proof}

\begin{lemma}[Open-region rank from geometric simple connectedness]%
    \label{thm:peps_torus_simply_connected_region_rank}
    \lean{TNLean.PEPS.IsGIsometric.rank_torusClosureRegionMatrix_of_isSimplyConnected}
    \leanok
    Let both torus circumferences be at least three and let $A$ be a
    native regular $G$-isometric four-leg tensor. Suppose the induced
    native graph on a finite region $R$ is connected and the
    closed-cell realization of $R$ in the real torus is simply connected.
    For a commuting closure pair $g,h$ and a numbering of its $b>0$
    crossing bonds, the twisted open-region coefficient matrix satisfies
    \begin{align}
        \operatorname{rank}M_R^{g,h} & =|G|^{b-1}.
    \end{align}
    No arithmetic lift, residual-flatness condition, or block Gram
    identity is assumed. Geometric simple connectedness and native
    graph connectivity are separate assumptions: closed cells may
    meet at diagonal points. This is an auxiliary open-region rank
    consequence of the boundary calculation in
    \cite[Theorem~6.9]{Schuch2010PEPS}, not the full disk entropy theorem
    or a coherent-superposition result
    \cite{gap:rmp_peps_quantum_double_g_isometry}.
\end{lemma}

\begin{proof}\leanok
    \uses{thm:peps_torus_region_realization, thm:peps_torus_region_integer_lift_gram}
    Graph connectivity gives a region vertex, and simple connectedness
    of the cell realization gives an integer lift preserving
    every internal native unit step by
    Lemma~\ref{thm:peps_torus_region_realization}. Applying the
    integer-lifted region rank theorem,
    Lemma~\ref{thm:peps_torus_region_integer_lift_gram}, to this lift completes the proof.
\end{proof}

\begin{lemma}[Canonical coordinates of a lifted torus region]%
    \label{thm:peps_torus_region_integer_lift_coordinates}
    \lean{TNLean.PEPS.regularProjectorTwistedRegionMatrix_coordinates_of_torusRegionIntegerLift,
        TNLean.PEPS.regularProjectorTwistedRegionMatrix_exists_normalized_ancilla_of_torusRegionIntegerLift}
    \leanok
    \uses{thm:peps_torus_region_integer_lift_gauge}
    Suppose a native torus region has the integer lift described in
    Lemma~\ref{thm:peps_torus_region_integer_lift_gauge}. Choose a spanning
    tree and root, and let $k^{g,h}$ be its normalized lift gauge. For
    commuting closure labels $g,h$, the canonical tensor satisfies
    \begin{align}
        B_R^{g,h}(y,a,r,z;\theta)
         & =B_R(y,a,r,z;F_{g,h}(\theta)),
    \end{align}
    where $F_{g,h}$ is the native transformed crossing configuration.
    In particular, there are a positive constant $\gamma$ and one
    normalized vector $\phi(a,r,z)$, independent of $g,h$, such that
    \begin{align}
        B_R^{g,h}(y,a,r,z;\theta)
         & =\gamma\Pi_{\partial R}(y,F_{g,h}(\theta))\phi(a,r,z).
    \end{align}
    This is an auxiliary lifted-region form of the blocking and
    boundary-coordinate calculations in \cite[\S6.2--6.3]{Schuch2010PEPS}.
    It retains the sector-dependent boundary map. The covering-lift
    and complementary-loop theorems provide the geometric choices for
    simply connected regions with torus periods at least three.
\end{lemma}

\begin{proof}\leanok
    \uses{thm:peps_regular_twisted_projector_coordinates, thm:peps_regular_region_projector_coordinates}
    Every cycle residual is the identity by
    Lemma~\ref{thm:peps_torus_region_integer_lift_gauge}. The exact twisted
    coefficient formula therefore reduces to the untwisted coefficient
    formula at the transformed crossing configuration. Applying the
    normalized factorization of
    Lemma~\ref{thm:peps_regular_region_projector_coordinates} to this
    configuration yields a constant and ancillary vector depending only on
    the region, the group, and the chosen tree and root.
\end{proof}
